% Options for packages loaded elsewhere
\PassOptionsToPackage{unicode}{hyperref}
\PassOptionsToPackage{hyphens}{url}
%
\documentclass[
  11pt,
]{article}
\usepackage{amsmath,amssymb}
\usepackage{iftex}
\ifPDFTeX
  \usepackage[T1]{fontenc}
  \usepackage[utf8]{inputenc}
  \usepackage{textcomp} % provide euro and other symbols
\else % if luatex or xetex
  \usepackage{unicode-math} % this also loads fontspec
  \defaultfontfeatures{Scale=MatchLowercase}
  \defaultfontfeatures[\rmfamily]{Ligatures=TeX,Scale=1}
\fi
\usepackage{lmodern}
\ifPDFTeX\else
  % xetex/luatex font selection
\fi
% Use upquote if available, for straight quotes in verbatim environments
\IfFileExists{upquote.sty}{\usepackage{upquote}}{}
\IfFileExists{microtype.sty}{% use microtype if available
  \usepackage[]{microtype}
  \UseMicrotypeSet[protrusion]{basicmath} % disable protrusion for tt fonts
}{}
\makeatletter
\@ifundefined{KOMAClassName}{% if non-KOMA class
  \IfFileExists{parskip.sty}{%
    \usepackage{parskip}
  }{% else
    \setlength{\parindent}{0pt}
    \setlength{\parskip}{6pt plus 2pt minus 1pt}}
}{% if KOMA class
  \KOMAoptions{parskip=half}}
\makeatother
\usepackage{xcolor}
\usepackage[letterpaper,margin=2.5cm]{geometry}
\usepackage{graphicx}
\makeatletter
\def\maxwidth{\ifdim\Gin@nat@width>\linewidth\linewidth\else\Gin@nat@width\fi}
\def\maxheight{\ifdim\Gin@nat@height>\textheight\textheight\else\Gin@nat@height\fi}
\makeatother
% Scale images if necessary, so that they will not overflow the page
% margins by default, and it is still possible to overwrite the defaults
% using explicit options in \includegraphics[width, height, ...]{}
\setkeys{Gin}{width=\maxwidth,height=\maxheight,keepaspectratio}
% Set default figure placement to htbp
\makeatletter
\def\fps@figure{htbp}
\makeatother
\setlength{\emergencystretch}{3em} % prevent overfull lines
\providecommand{\tightlist}{%
  \setlength{\itemsep}{0pt}\setlength{\parskip}{0pt}}
\setcounter{secnumdepth}{-\maxdimen} % remove section numbering
\ifLuaTeX
\usepackage[bidi=basic]{babel}
\else
\usepackage[bidi=default]{babel}
\fi
\babelprovide[main,import]{spanish}
% get rid of language-specific shorthands (see #6817):
\let\LanguageShortHands\languageshorthands
\def\languageshorthands#1{}
% Ajustes de calidad editorial para los boletines IES.
% Se incluyen con: pandoc --include-in-header=assets/editorial.tex
% Nota: pandoc ya carga `microtype` automáticamente si está disponible,
% por eso aquí solo se afinan los pies de figura.
\usepackage{caption}     % control de los pies de figura
\captionsetup{font=small, labelfont=bf}
\ifLuaTeX
  \usepackage{selnolig}  % disable illegal ligatures
\fi
\IfFileExists{bookmark.sty}{\usepackage{bookmark}}{\usepackage{hyperref}}
\IfFileExists{xurl.sty}{\usepackage{xurl}}{} % add URL line breaks if available
\urlstyle{same}
\hypersetup{
  pdflang={es},
  hidelinks,
  pdfcreator={LaTeX via pandoc}}

\author{}
\date{}

\begin{document}

\hypertarget{termodinuxe1mica-financiera-05-boletuxedn-toc-insights}{%
\section{TERMODINÁMICA FINANCIERA 05: BOLETÍN TOC
INSIGHTS}\label{termodinuxe1mica-financiera-05-boletuxedn-toc-insights}}

\textbf{Análisis Sistémico de PEMEX: La Propuesta de Escisión bajo la
Teoría de Restricciones}

\href{https://es-us.noticias.yahoo.com/proponen-escindir-pemex-negocios-deficitarios-050356273.html}{Escision
de pemex}

\begin{center}\rule{0.5\linewidth}{0.5pt}\end{center}

\hypertarget{titular-y-resumen-ejecutivo}{%
\subsubsection{1. Titular y Resumen
Ejecutivo}\label{titular-y-resumen-ejecutivo}}

La consultora Welligence propuso escindir los negocios deficitarios de
Petróleos Mexicanos (PEMEX), evaluando a la empresa estatal bajo los
criterios de rentabilidad de la contabilidad de costos privada. Desde la
Teoría de Restricciones (TOC), desmembrar la infraestructura bajo la
premisa de aislar pérdidas ignora el rol sistémico de la petrolera como
amortiguador macroeconómico y escudo energético nacional, confundiendo
la optimización contable local con la maximización del \emph{Throughput}
y el bienestar social.

\begin{center}\rule{0.5\linewidth}{0.5pt}\end{center}

\hypertarget{diagnuxf3stico-del-sistema-seguxfan-toc}{%
\subsubsection{2. Diagnóstico del Sistema según
TOC}\label{diagnuxf3stico-del-sistema-seguxfan-toc}}

\begin{itemize}
\item
  \textbf{Restricción del Sistema (Cuello de Botella):} El cuello de
  botella operativo y financiero no proviene de la diversidad de
  unidades de negocio, sino de una restricción dual: la insuficiencia de
  capital destinado a la perforación y exploración rentable (con una
  caída drástica en la perforación de pozos), estrangulado por una carga
  fiscal asfixiante y el desvío de recursos hacia pasivos, deudas y
  refinación no sincronizada.
\item
  \textbf{Parámetros Clave del Sistema:}
\item
  \textbf{Throughput (\(T\)):} El valor energético y la estabilidad
  macroeconómica provistos al país para amortiguar shocks externos de
  precios (como la contención de precios de gasolinas frente a tensiones
  geopolíticas globales).
\item
  \textbf{Inversión / Inventario (\(I\)):} Activos físicos estratégicos
  (campos de producción, refinerías, red logística) y la deuda
  financiera acumulada.
\item
  \textbf{Gasto de Operación (\(OE\)):} Costos asociados al
  mantenimiento de la estructura corporativa, nómina y operación de los
  complejos.
\item
  \textbf{Representación en d2lang de la Nube de Evaporación del
  Conflicto:}
\end{itemize}

\begin{figure}
\centering
\includegraphics{Diagramas/figura-1.pdf}
\caption{GARANTIZAR BIENESTAR SOCIAL Y SOBERANÍA ENERGÉTICA}
\end{figure}

\begin{center}\rule{0.5\linewidth}{0.5pt}\end{center}

\hypertarget{dicotomuxeda-clave-1-contabilidad-de-costos-tradicional-vs.-contabilidad-del-throughput-ta}{%
\subsubsection{3. Dicotomía Clave 1: Contabilidad de Costos Tradicional
vs.~Contabilidad del Throughput
(TA)}\label{dicotomuxeda-clave-1-contabilidad-de-costos-tradicional-vs.-contabilidad-del-throughput-ta}}

\begin{itemize}
\tightlist
\item
  \textbf{Visión de la Contabilidad de Costos Tradicional:} Analiza de
  forma aislada cada segmento de PEMEX (exploración, refinación,
  petroquímica). Si una unidad reporta márgenes negativos o absorbe
  flujos, exige su escisión, venta o cierre para ``mejorar el EBITDA
  corporativo'' y reducir costos por componente u operación local.
\item
  \textbf{Visión de la Contabilidad del Throughput (TA):} Advierte que
  la reducción de costos locales (\(OE\)) o la eliminación de activos de
  soporte bajo criterios fragmentados puede colapsar el flujo global del
  sistema. Un segmento con déficit contable interno puede operar como el
  amortiguador indispensable que protege el \emph{Throughput}
  macroeconómico y evita que la volatilidad de precios colapse el
  aparato productivo nacional.
\end{itemize}

\begin{center}\rule{0.5\linewidth}{0.5pt}\end{center}

\hypertarget{dicotomuxeda-clave-2-teoruxeda-monetaria-moderna-mmt-vs.-teoruxeda-de-restricciones-toc}{%
\subsubsection{4. Dicotomía Clave 2: Teoría Monetaria Moderna (MMT)
vs.~Teoría de Restricciones
(TOC)}\label{dicotomuxeda-clave-2-teoruxeda-monetaria-moderna-mmt-vs.-teoruxeda-de-restricciones-toc}}

\begin{itemize}
\tightlist
\item
  \textbf{Visión MMT:} Al tratarse de un Estado con soberanía monetaria,
  las restricciones financieras nominales son secundarias; los apoyos
  fiscales, la absorción de deudas o los déficits operativos en empresas
  públicas no están limitados por ingresos tributarios previos, sino por
  los recursos reales disponibles y el riesgo inflacionario.
\item
  \textbf{Visión TOC:} El dinero es simplemente un medio facilitador; la
  restricción fundamental es física y operativa. Incluso con soberanía
  monetaria, si la capacidad de extracción y refinación (la verdadera
  restricción del sistema energético) está estrangulada por una
  asignación deficiente de recursos hacia gastos no productivos en lugar
  de pozos de alta rentabilidad, la inyección de liquidez o el
  endeudamiento solo elevan el inventario financiero (\(I\)) y generan
  ruido inflacionario sin acelerar el flujo físico de bienestar.
\end{itemize}

\begin{center}\rule{0.5\linewidth}{0.5pt}\end{center}

\hypertarget{conclusiuxf3n-y-prescripciuxf3n-estratuxe9gica-5-pasos-de-enfoque-de-goldratt}{%
\subsubsection{5. Conclusión y Prescripción Estratégica (5 Pasos de
Enfoque de
Goldratt)}\label{conclusiuxf3n-y-prescripciuxf3n-estratuxe9gica-5-pasos-de-enfoque-de-goldratt}}

\begin{enumerate}
\def\labelenumi{\arabic{enumi}.}
\tightlist
\item
  \textbf{Identificar la Restricción:} Reconocer que el cuello de
  botella no es la estructura integrada de la empresa pública, sino la
  distorsión en la asignación de recursos (fondos desviados al
  saneamiento de pasivos e impuestos en lugar de la perforación de
  campos altamente productivos).
\item
  \textbf{Explotar la Restricción:} Maximizar la eficiencia operativa de
  los campos y activos actuales que concentran la mayor parte de la
  producción para asegurar el suministro energético interno frente a la
  volatilidad geopolítica.
\item
  \textbf{Subordinar el Resto del Sistema:} Alinear la política fiscal
  gubernamental para detener la extracción excesiva de caja (vía
  impuestos confiscatorios o uso como caja chica), subordinando las
  metas de recaudación de corto plazo a la estabilidad del flujo
  energético nacional.
\item
  \textbf{Elevar la Restricción:} Inyectar capital de inversión (\(I\))
  de forma directa y prioritaria en la perforación, tecnología y
  mantenimiento de los campos rentables, ampliando la capacidad de
  producción del amortiguador macroeconómico.
\item
  \textbf{Prevenir la Inercia:} Evitar la aplicación ciega de dogmas
  corporativos privados que dictan la fragmentación de empresas
  estatales bajo métricas de contabilidad de costos sin evaluar el
  impacto sistémico sobre la estabilidad del país y el bienestar social.
\end{enumerate}

\end{document}
